\section{Square grids}\label{app:square}

This appendix proves Theorem~\ref{thm:square}. Throughout, $n\ge2$, the grid is $n\times n$, $\mathcal R$ is the set of its rectangles, tilted ones included, and $L=1+\ln n$ with the natural logarithm. For a point $v$ we write $d_v$ for the number of rectangles that have $v$ as a corner. By Lemma~\ref{lem:threecorners}, two different rectangles share at most two corners, so the symmetric difference of their corner sets has four, six, or eight points.

\subsection{Pairs with two common corners}

\begin{lemma}[two common corners]\label{lem:compress}
Let $R_1\ne R_2$ be rectangles with exactly two common corners $p$ and $q$. If $\{p,q\}$ is a side of both rectangles, or a diagonal of both, then $C(R_1)\triangle C(R_2)$ is the corner set of a rectangle.
\end{lemma}

\begin{proof}
Suppose $\{p,q\}$ is a side of both. The other corners are $p+u,q+u$ for $R_1$ and $p+v,q+v$ for $R_2$, where $u\ne v$ are nonzero vectors perpendicular to $q-p$. The four points $p+u$, $q+u$, $p+v$, $q+v$ are distinct, since $p+u=q+v$ would make $v-u=p-q$ a nonzero vector perpendicular to itself. None of them is $p$ or $q$, since the common corners are exactly $p$ and $q$. They form a parallelogram with sides $q-p$ and $v-u$, which are nonzero and perpendicular, so they are the corners of a rectangle.

Suppose $\{p,q\}$ is a diagonal of both. Both rectangles have center $o=(p+q)/2$, and their other corners are $o\pm u$ for $R_1$ and $o\pm v$ for $R_2$ with $|u|=|v|=|p-q|/2$ and $v\ne\pm u$. The segments from $o-u$ to $o+u$ and from $o-v$ to $o+v$ have the same midpoint and the same length. So they are the diagonals of a rectangle with corners $o\pm u$, $o\pm v$, four distinct grid points.
\end{proof}

\subsection{Lattice points on circles}

Let $r_2(t)$ be the number of pairs $(a,b)$ of integers with $a^2+b^2=t$, and let $\tau(t)$ be the number of positive divisors of $t$.

\begin{lemma}[divisor bound]\label{lem:divisor}
For every $t\ge1$, $r_2(t)\le4\tau(t)$. For every $\varepsilon>0$ there is a constant $K_\varepsilon$ such that $\tau(t)\le K_\varepsilon t^\varepsilon$ for all $t\ge1$.
\end{lemma}

\begin{proof}
Both parts are classical. Hardy and Wright \cite[Section~16.9]{hardywright2008} prove that $r_2(t)=4\bigl(d_1(t)-d_3(t)\bigr)$, where $d_i(t)$ is the number of divisors of $t$ congruent to $i$ modulo $4$, so $r_2(t)\le4d_1(t)\le4\tau(t)$. The bound $\tau(t)=O(t^\varepsilon)$ for every fixed $\varepsilon>0$ is proved in Section~18.1 of the same book.
\end{proof}

\subsection{Counting rectangles and obstructions}

Let $\mathcal A$ be the family of four-point sets consisting of all corner sets $C(R)$ with $R\in\mathcal R$, together with the mixed differences. A \emph{mixed difference} is a set $C(R_1)\triangle C(R_2)$ for a pair $R_1\ne R_2$ with exactly two common corners that form a side of one rectangle and a diagonal of the other, the case that Lemma~\ref{lem:compress} leaves out. Let $\mathcal B$ be the family of all six-point sets $C(R_1)\triangle C(R_2)$ for pairs $R_1\ne R_2$ with exactly one common corner.

\begin{lemma}[rectangle and obstruction counts]\label{lem:counts}
Let $n\ge2$ and $L=1+\ln n$.
\begin{enumerate}[(a)]
\item For every point $v$, $d_v\le16n^2L$, and $|\mathcal R|\le4n^4L$.
\item For every $\varepsilon>0$, $|\mathcal A|\le A_\varepsilon n^{4+2\varepsilon}L$ with $A_\varepsilon=4(1+16\cdot2^\varepsilon K_\varepsilon)$, and $|\mathcal B|\le128n^6L^2$.
\item If a blackout $S$ of the $n\times n$ grid contains no member of $\mathcal A$ and no member of $\mathcal B$, then $S$ is valid.
\end{enumerate}
\end{lemma}

\begin{proof}
(a) Let $R$ be a rectangle with corner $v$. Its two sides from $v$ are perpendicular nonzero integer vectors. Write the first as $k(a,b)$ with $k\ge1$ and $\gcd(a,b)=1$. The second is perpendicular to $(a,b)$, so it is $(-bl,al)$ for a real $l$, and $l$ is a nonzero integer because the vector is integral and $\gcd(a,b)=1$. Put $d=\max(|a|,|b|)\ge1$. There are $8d$ integer vectors $(a,b)$ with $\max(|a|,|b|)=d$, since they lie on the boundary of the square $[-d,d]^2$. Both sides fit in the grid, so $kd\le n-1$ and $|l|d\le n-1$, which leaves at most $(n-1)/d$ values of $k$ and $2(n-1)/d$ values of $l$. The ordered pair of sides determines $R$, so
\[
d_v\le\sum_{d=1}^{n-1}8d\cdot\frac{n-1}{d}\cdot\frac{2(n-1)}{d}\le16n^2\sum_{d=1}^{n-1}\frac1d\le16n^2L.
\]
Each rectangle has four corners, so $4|\mathcal R|=\sum_vd_v\le n^2\cdot16n^2L$.

(b) Fix $R_1\in\mathcal R$ and one of its four sides $\{p,q\}$. A rectangle $R_2$ that has $\{p,q\}$ as a diagonal has center $(p+q)/2$ and a second diagonal of the same length, joining two grid points $p',q'$. The vector $p'-q'$ is an integer vector with $|p'-q'|^2=|p-q|^2=t$, where $t\le2(n-1)^2<2n^2$, and it determines $p'$ and $q'$. So there are at most $r_2(t)\le4K_\varepsilon(2n^2)^\varepsilon$ such $R_2$ by Lemma~\ref{lem:divisor}. Every mixed difference arises from some $R_1$, side, and $R_2$ in this way, so
\[
|\mathcal A|\le|\mathcal R|\bigl(1+4\cdot4K_\varepsilon2^\varepsilon n^{2\varepsilon}\bigr)\le4n^4L\bigl(1+16\cdot2^\varepsilon K_\varepsilon\bigr)n^{2\varepsilon}=A_\varepsilon n^{4+2\varepsilon}L .
\]
A member of $\mathcal B$ comes from an unordered pair of rectangles with a common corner $v$, so
\[
|\mathcal B|\le\sum_v\binom{d_v}2\le\frac{n^2}2\bigl(16n^2L\bigr)^2=128n^6L^2 .
\]

(c) Suppose $S$ is not valid, so some $R_1\ne R_2$ have the same presentation, and $C(R_1)\triangle C(R_2)\subseteq S$ by Lemma~\ref{lem:diff}. If $R_1$ and $R_2$ have two common corners, the symmetric difference is the corner set of a rectangle by Lemma~\ref{lem:compress} or a mixed difference, and in both cases it is a member of $\mathcal A$. If they have one common corner, it is a member of $\mathcal B$. If they have none, then $C(R_1)\subseteq S$ is a member of $\mathcal A$. In each case $S$ contains a member of $\mathcal A\cup\mathcal B$.
\end{proof}

Avoiding every member of $\mathcal A$ forbids every fully black rectangle, while the game allows one. This is harmless for a lower bound, and it removes the eight-point symmetric differences from the count.

\subsection{The lower bound}

\begin{theorem}[random blackouts with deletions]\label{thm:alteration}
Let $\alpha=512^{-1/5}$. There is $n_0$ such that for all $n\ge n_0$,
\[
f(n,n)\ge\frac\alpha2\,\frac{n^{6/5}}{(1+\ln n)^{2/5}} .
\]
In particular $f(n,n)-2n\to\infty$.
\end{theorem}

\begin{proof}
Put $\theta=\alpha\,n^{-4/5}L^{-2/5}$, which lies in $(0,1)$ since $\alpha<1$, and select each point of the grid independently with probability $\theta$. Let $X$ be the number of selected points. Let $Y$ be the number of members of $\mathcal A$ all of whose four points are selected, and $Z$ the number of members of $\mathcal B$ all of whose six points are selected. By linearity of expectation,
\begin{equation}\label{eq:expect}
\mathbb E(X-Y-Z)=n^2\theta-|\mathcal A|\theta^4-|\mathcal B|\theta^6 .
\end{equation}
We compare the two losses with $n^2\theta=\alpha n^{6/5}L^{-2/5}$. By Lemma~\ref{lem:counts}(b),
\[
\frac{|\mathcal B|\theta^6}{n^2\theta}\le\frac{128n^6L^2\,\alpha^5n^{-4}L^{-2}}{n^2}=128\alpha^5=\frac14 ,
\]
by the choice of $\alpha$. With $\varepsilon=1/10$ and $A_0=A_{1/10}$, Lemma~\ref{lem:counts}(b) gives $|\mathcal A|\le A_0n^{21/5}L$, so
\[
\frac{|\mathcal A|\theta^4}{n^2\theta}\le\frac{A_0n^{21/5}L\,\alpha^3n^{-12/5}L^{-6/5}}{n^2}=A_0\alpha^3\,n^{-1/5}L^{-1/5}\le\frac14
\]
as soon as $n\ge n_0=(4A_0\alpha^3)^5$, since $L\ge1$. For such $n$, \eqref{eq:expect} gives $\mathbb E(X-Y-Z)\ge n^2\theta/2$, so some selection $S_0$ has $X-Y-Z\ge n^2\theta/2$.

Fix this $S_0$. Go through the members of $\mathcal A\cup\mathcal B$ that lie inside $S_0$, and for each one that is still contained in the current set delete one of its points. This deletes at most $Y+Z$ points and leaves a set $S\subseteq S_0$ with $|S|\ge X-Y-Z\ge n^2\theta/2$. No member of $\mathcal A\cup\mathcal B$ is contained in $S$. Indeed a member contained in $S$ is contained in $S_0$, so it was visited, and at that moment it lost a point or had already lost one, and the families do not depend on the deletions. By Lemma~\ref{lem:counts}(c), $S$ is a valid blackout, so $f(n,n)\ge n^2\theta/2=\frac\alpha2n^{6/5}L^{-2/5}$. Finally $n^{6/5}L^{-2/5}/n=n^{1/5}L^{-2/5}\to\infty$, which gives the last claim.
\end{proof}

Every $\varepsilon<1/5$ makes the first loss ratio $O(n^{2\varepsilon-2/5})$, which tends to $0$, and $\varepsilon=1/10$ is one convenient value. The method is the alteration method of Alon and Spencer \cite[Section~3.2]{alonspencer2016}. The six-point family is the kind of obstruction for which Spencer's bound on independent sets in uniform hypergraphs \cite{spencer1972}, in the form quoted by Conlon, Fox, and Sudakov \cite[Section~2]{conlonfoxsudakov2016}, gives the exponent $6/5$, and Lemmas~\ref{lem:compress} and~\ref{lem:counts} keep the four-point family small enough.

\subsection{The upper bound}

\begin{proposition}[elementary upper bound]\label{prop:upper}
For every $n\ge3$, $f(n,n)\le2n^{4/3}+2n$.
\end{proposition}

\begin{proof}
Let $S$ be a valid blackout of the $n\times n$ grid with $e=|S|$ points, and let $G=G_S$, a bipartite graph on $2n$ vertices with $e$ edges. Since $n\ge3$, Lemma~\ref{lem:girth} says that $G$ has no cycle of length four or six, and as a bipartite graph it has no odd cycle, so $G$ has no cycle of length at most seven. If $e=0$ there is nothing to prove, so let $e>0$ and put $t=e/(2n)$.

\emph{Step 1. A subgraph of large minimum degree.} Delete vertices one at a time, each time a vertex whose current degree is less than $t$, until no such vertex remains. If every vertex were deleted, every edge would be removed at the deletion of the first of its ends, a vertex of current degree less than $t$, so $e<2n\cdot t=e$. So a nonempty subgraph $G'$ remains, with minimum degree $\delta\ge t>0$. Let $uv$ be an edge of $G'$.

\emph{Step 2. The vertices near an edge.} Distances below are in $G'$. Since $u$ and $v$ are adjacent and $G'$ is bipartite, every vertex $x$ in the component of $u$ has distances to $u$ and to $v$ that differ by exactly one. For $0\le i\le3$ let $U_i$ be the set of vertices at distance $i$ from $u$ and $i+1$ from $v$, and $V_i$ the set at distance $i$ from $v$ and $i+1$ from $u$. These eight sets are pairwise disjoint, and $U_0=\{u\}$, $V_0=\{v\}$. We claim that $|U_{i+1}|\ge(\delta-1)|U_i|$ for $i=0,1,2$.

Let $x\in U_i$ with $i\le2$. At most one neighbor of $x$ is at distance $i-1$ from $u$, since two such neighbors would give two different paths of length $i$ from $u$ to $x$, which contain a cycle of length at most $2i\le4$. For $i=0$ the one excluded neighbor is $v$ itself. No neighbor of $x$ is at distance $i$ from $u$, by parity. So at least $\delta-1$ neighbors $y$ of $x$ are at distance $i+1$ from $u$, and for such a $y$ the distance to $v$ is $i$ or $i+2$. Suppose it were $i$. Then the path from $u$ to $x$ of length $i$ followed by the edge $xy$, and the edge $uv$ followed by a path from $v$ to $y$ of length $i$, are two paths of length $i+1$ from $u$ to $y$. They are different, because the second passes through $v$ and the first does not. Indeed a vertex of a shortest path from $u$ to $x$ is at distance at most $i$ from $x$, and $y\ne v$ when $i=0$. So they would contain a cycle of length at most $2i+2\le6$. Hence $y\in U_{i+1}$. Finally two different vertices $x,x'\in U_i$ have no common neighbor $y\in U_{i+1}$, since again this would give two different paths of length $i+1$ from $u$ to $y$. This proves the claim, and the same holds for the sets $V_i$.

\emph{Step 3. The count.} The eight disjoint sets lie in $G'$, which has at most $2n$ vertices, so $2(1+(\delta-1)+(\delta-1)^2+(\delta-1)^3)\le2n$. Hence $(\delta-1)^3\le n$, and $e=2nt\le2n\delta\le2n(1+n^{1/3})=2n+2n^{4/3}$.
\end{proof}

\begin{proof}[Proof of Theorem~\ref{thm:square}]
Part (a) is Theorem~\ref{thm:alteration} with $c=\alpha/2=512^{-1/5}/2$, and part (b) is Proposition~\ref{prop:upper}.
\end{proof}
